\chapter{Этапы выполнения}
\label{cha:stages}

Для выполнения поставленных задач были сформулированы следующие технологические требования к инструментарию:
\begin{itemize}
    \item Поддержка реляционной модели данных и стандарта SQL в СУБД PostgreSQL~\cite{doc:postgres,doc:sql};
    \item Поддержка сетевого протокола передачи гипертекста HTTP~\cite{wiki:http};
    \item Возможность отрисовки форм, кнопок и динамических таблиц в веб-интерфейсе.
\end{itemize}

Для реализации клиентского интерфейса был выбран язык разметки HTML~\cite{wiki:html} совместно с библиотекой стилей Pico CSS~\cite{wiki:pico.css}, обеспечивающей лаконичное визуальное оформление без необходимости написания сложного клиентского JavaScript-кода.

В качестве языка серверной части выбран Golang версии 1.27~\cite{wiki:go}. Для взаимодействия с СУБД PostgreSQL задействован драйвер и пул соединений \texttt{github.com/jackc/pgx}~\cite{git:pgx}.

Сервер функционирует по принципу Server-Side Rendering (SSR). Все страницы собираются на стороне сервера с помощью встроенного шаблонизатора \texttt{html/template} и возвращаются браузеру в виде готового HTML-документа.
